\chapter{Аналитическая часть}
\section{Расстояние Левенштейна}
Расстояние Левенштейна~---~минимальное количество редакторских операций вставки, удаления и замены одного символа, необходимых для преобразования одной строки к другой.

Вводятся операции слудющих типов:
\begin{enumerate}
	\item I (англ. insert)~---~операция вставки символа;
	\item D (англ. delete)~---~операция удаления символа;
	\item R (англ. replace)~---~операция замены символа.
\end{enumerate}

Стоимость каждой их этих операций принимается равной 1.

Вводятся также следующие обозначения:
\begin{itemize}[label=---]
	\item $\lambda$~---~пустой символ, не входящий ни в одну из рассматриваемых строк;
	\item D(a, b)~---~расстояние Левенштейна от двух строк a и b;
	\item M (англ. match)~---~совпадение символов.
\end{itemize}

Тогда:
\begin{enumerate}
	\item $D(a, b) = 1$, где $a \neq b$;
	\item $D(a, b) = 0$, где $a = b$;
	\item $D(\lambda, b) = 1$;
	\item $D(a, \lambda) = 1$.
\end{enumerate}

Существует проблема взаимного выравнивания символов строк. Эта проблема решается введением рекуррентной формулы. 

Реккурентная формула~---~формула, выражающая следующие члены рода или последовательности через предыдущие.

Пусть даны строки $S_1$, $S_2$, а также
\begin{itemize}[label=---]
	\item $L_1$ --- длина $S_1$;
	\item $L_2$ --- длина $S_2$;
	\item $S_1[1...i]$ --- подстрока $S_1$ длиной $i$, начиная с первого символа;
	\item $S_2[1...j]$ --- подстрока $S_2$ длиной $j$, начиная с первого символа;
\end{itemize}

Расстояние Левенштейна между двумя строками $S_1$ и $S_2$ длиной $L_1$ и $L_2$ соответственно рассчитывается по рекуррентной формуле (\ref{eq:lev_req}):
\small
\begin{equation}
	\label{eq:lev_req}
	D(S_1[1..i],S_2[1..j]) = 
	\begin{cases}
		\begin{array}{ll}
			\max(i,j), & \textrm{$\mbox{если }\min(i,j) = 0,$}\\
			\min(D(S_1[1..i], S_2[1.. j - 1])+1,\\
			\qquad D(S_1[1..i - 1], S_2[1..j]) + 1, &\textrm{$\mbox{иначе},$}\\
			\qquad D(S_1[1..i - 1], S_2[1..j - 1]) +\\
			\qquad + not\_eq(S_1[i],S_2[j])),
		\end{array}
	\end{cases}
\end{equation}
где сравнение символов строк $S_1$ и $S_2$ рассчитывается как:
\begin{equation}
	\label{eq:not\_eq}
	not\_eq(a, b) = \begin{cases}
		0, &\text{если a = b,}\\
		1, &\text{иначе.}
	\end{cases}
\end{equation}

\normalsize
Обозначим результат подсчета расстояния Левенштейна для подстрок $S_1[1..i],S_2[1..j]$ как 
$Lev(S_1[1..i],S_2[1..j])$.

\section{Расстояние Дамерау---Левенштейна}

В отличии от расстояния Левенштейна, в расстоянии Дамерау---Левенштейна вводится еще одна операция --- транспонирования $T$ (англ. transposition). Эта операция применима,  если $S_1[i] = S_2[j - 1]$ и $S_1[i - 1] = S_2[j]$. Рекуррентная формула данной метрики выглядит следующим образом:
\small
\begin{equation}
	\label{eq:dam_lev_req}
	D(S_1[1..i],S_2[1..j]) = 
	\begin{cases}
		min \begin{cases}
			Lev(S_1[1..i],S_2[1..j]),\\
			D(S_1[1..i - 2],S_2[1..j - 2]) + 1, \\
		\end{cases}
		& \begin{aligned}
			& \text{если i > 1, j > 1}, \\
			& S_1[i] = S_2[j - 1], \\
			& S_1[i - 1] = S_2[j], \\
		\end{aligned}\\
		Lev(S_1[1..i],S_2[1..j])
		& \text{иначе.}
	\end{cases}
\end{equation}
